\documentclass[12pt,a4paper]{article}
\usepackage{graphicx}
\usepackage{amsmath}
\usepackage{bm}
\usepackage{interval}
\usepackage{amssymb}
\usepackage{float}
\usepackage{subfigure}
\usepackage{titlesec}
% \usepackage{algorithm}
\usepackage[linesnumbered,boxed]{algorithm2e}
% \usepackage{algorithmic}
\usepackage[noend]{algpseudocode}
\usepackage{epsfig}
\usepackage[letterpaper,top=2cm,bottom=2cm,left=3cm,right=3cm,marginparwidth=1.75cm]{geometry}
\usepackage[colorlinks=true, allcolors=blue]{hyperref}

% Inputs, Initialize settings
\algnewcommand{\Inputs}[1]{%
  \State \textbf{Inputs:}
  \Statex \hspace*{\algorithmicindent}\parbox[t]{.8\linewidth}{\raggedright #1}
}
\algnewcommand{\Initialize}[1]{%
  \State \textbf{Initialize:}
  \Statex \hspace*{\algorithmicindent}\parbox[t]{.8\linewidth}{\raggedright #1}
}

% Change section font size
\titleformat{\section}  % which section command to format
  {\fontsize{14}{16}\bfseries} % format for whole line
  {\thesection} % how to show number
  {1em} % space between number and text
  {} % formatting for just the text
  [] % formatting for after the text

\title{DS Written Homework 4}
\author{Student ID: b10201054}

\begin{document}
\maketitle

\section{AA Tree}



\section{Reb-Black Tree}    

% \begin{enumerate}
%     \item[(a)] 
%         \begin{figure}[H]
%             \centering
%             \includegraphics[width=0.7\linewidth]{Splaytree1.png}
%             \caption{Splaytree(a)}
%         \end{figure}
        
%     \item[(b)]
%         \begin{figure}[H]
%             \centering
%             \includegraphics[width=0.7\linewidth]{Splaytree2.png}
%             \caption{Splaytree(b)}
%         \end{figure}
    
%     \item[(c)]
%         \begin{figure}[H]
%             \centering
%             \includegraphics[width=0.7\linewidth]{Splaytree3.png}
%             \caption{Splaytree(c)}
%         \end{figure}
    
% \end{enumerate}

\section{The Height of Red-Black Tree}    

\begin{enumerate}
    \item[(a)] 
        
    \item[(b)]
    
    \item[(c)]
     
    \item[(d)] 
    
    \item[(e)]

\end{enumerate}

\section{2-4 Tree}

% \begin{enumerate}
%     \item[(a)] 
%         \begin{figure}[H]
%             \centering
%             \includegraphics[width=0.7\linewidth]{Splaytree1.png}
%             \caption{Splaytree(a)}
%         \end{figure}
        
%     \item[(b)]
%         \begin{figure}[H]
%             \centering
%             \includegraphics[width=0.7\linewidth]{Splaytree2.png}
%             \caption{Splaytree(b)}
%         \end{figure}
    
% \end{enumerate}

\section{Possible Red-Black Tree}

\begin{enumerate}
    \item[(a)] 
        
    \item[(b)]
    
    \item[(c)]

\end{enumerate}

\section{Hash}

\begin{enumerate}
    \item[(a)] 
        
    \item[(b)]
    
    \item[(c)]
     
    \item[(d)] 

\end{enumerate}

\section{Fill-in-the-blank questions}

\begin{enumerate}
    \item 
    \item 
    \item 
    \item 
    \item 
\end{enumerate}

\end{document}
